%% requires: \providecommand{\vfy}{\textsuperscript{[v]}}  (marks a value still to be verified)
% ============================================================
% GROUP 7: Potentials, validation (F7b) and robust ranking (F7)
% Place: Literature Review (section 1) and Methodology (section 2).
% Dropped from the agent draft: the NPV(K,s) formula, the chapter 3 equations (1) to (5)
% as written, and the Cabrera and Gianetti MAUP references (off topic).
% ============================================================

\section{Literature Review: Potentials, Spatial Scale and Robust Decision Rules}
\label{sec:g7-lit}

\subsection{Technical potential is a range that depends on premises}
Potential estimates follow a hierarchy from theoretical to technical, economic and feasible potential, and published values differ by orders of magnitude because of resolution, reanalysis data and assumptions on power density and exclusions \citep{angliviel2023}. The values available for the three study countries illustrate this (Table~\ref{tab:g7-pot}). For Brazilian onshore wind, the first national atlas reports 143~GW at 50~m \citep{amarante2001}, state atlases compiled in the 2050 plan sum to about 315~GW at 100~m and 598~GW at 150~m for twelve states \citep{epe2020pne}, and a global estimate with 5~MW/km$^2$ over 31\% of the territory gives 14.58~TW \citep{eurek2017}. For Indian wind, the national assessments rise with hub height from 302~GW at 100~m to 695.5~GW at 120~m and about 1164~GW at 150~m \citep{niwe_atlas}. These values are not comparable without harmonizing hub height, power density, exclusion rules and the capacity-factor cutoff. A difference between a modeled potential and a published one is therefore expected and is explained by premises, not treated as an error.

\begin{table}[ht]
\centering
\caption{Published technical potentials used as validation bands, with their premises.}
\label{tab:g7-pot}
\footnotesize
\begin{tabular}{p{1.5cm}p{1.1cm}p{2.0cm}p{5.8cm}p{2.0cm}}
\toprule
Country & Tech. & Potential & Main premises & Source \\
\midrule
Brazil & Wind & 143~GW & 50~m, 1~km grid, mean wind above 7~m/s (2001) & \citet{amarante2001} \\
Brazil & Wind & 315~GW; 598~GW & Twelve state atlases at 100~m; 150~m & \citet{epe2020pne}\vfy{} \\
Brazil & Wind & 14.58~TW (22.22~PWh/yr) & 3.5~MW, 90~m, 5~MW/km$^2$, 0.95 availability, 31\% of area available & \citet{eurek2017} \\
Brazil & PV & 307~GWp (506~TWh/yr) & Best irradiation band (6.0 to 6.2~kWh/m$^2$/day), anthropized areas only, 70~MWp/km$^2$, CF$_{\mathrm{DC}}$ 18.84\% & \citet{epe2018potencial,denholm2008} \\
Brazil & PV & 28\,519~GWp (43\,857~TWh/yr) & Eight irradiation bands, 407\,437~km$^2$; an area-based ceiling, not a deployable level & \citet{epe2018potencial} \\
India & PV & 748~GWp & 3\% of wasteland area (2014) & \citet{mnre2024pib748}\vfy{} \\
India & PV & about 3\,343~GWp & Ground-mounted, updated assessment (2025) & \citet{nise2025ground}\vfy{} \\
India & Wind & 302; 695.5; about 1\,164~GW & 100~m; 120~m; 150~m & \citet{niwe_atlas}\vfy{} \\
Portugal & PV, wind & not available & Area of low-sensitivity sites mapped; no totals found & \citet{simoes2023lneg}\vfy{} \\
\bottomrule
\end{tabular}
\end{table}

\subsection{Spatial aggregation}
Gridded results depend on the size and shape of the spatial unit. This modifiable areal unit problem has a scale effect and a zoning effect \citep{openshaw1984}. Simulation shows that units that are too coarse lose efficiency and misclassify, while units that are too fine relative to the data-generating process add attenuation bias \citep{avelino2016}. In a downscaling application the scale effect was larger than the zoning effect, and spatially clustered variables were less affected \citep{dare2020}. A formal criterion for the information lost by aggregating between scales has been proposed using a multiscale expansion \citep{bradley_cage}. For a 0.05$^{\circ}$ grid (about 5.5~km) these results motivate a sensitivity test of the ranking to coarser aggregation, not a single fixed resolution.

\subsection{Regret, satisficing and scenario discovery}
Under deep uncertainty, expected-value optimization is replaced by robustness criteria. Minimax regret chooses the option with the smallest worst-case loss relative to the best option in each future \citep{savage1951,lempert2006}. Satisficing or domain criteria count the share of futures in which performance meets a required level \citep{herman2015,mcphail2018}, and info-gap theory expresses the same idea as the size of the uncertainty horizon that can be tolerated \citep{benhaim2006}. Patient rule induction (PRIM) finds boxes in the space of uncertain factors where a failure outcome concentrates \citep{friedman1999,bryant2010}. Many-objective and adaptive variants are reviewed in the decision-making literature \citep{kasprzyk2013,hamarat2013,marchau2019}. In energy, minimax regret has been applied to greenhouse gas abatement strategies \citep{loulou1999}, regret and info-gap have been compared for climate policy \citep{hall2012}, and exploratory robust design has been applied to energy transitions \citep{hamarat2013}. Among the sources reviewed for this work, no peer-reviewed application combining regret, satisficing and PRIM for spatial siting of wind or PV on a grid was identified. This is a statement about the sources reviewed, not a claim of novelty.

Minimax regret has known limits. It can violate independence of irrelevant alternatives and can produce non-transitive rankings, and for a continuous, convex decision set with quasi-convex regret the optimum is determined by at most $n+1$ scenarios, where $n$ is the dimension of the decision. That reduction is shown not to hold for discrete choices \citep{andersonzachary}. % verify: reading of the theorem taken from the research agent's summary of the PDF
A ranking of grid cells is a discrete choice, so a small set of extreme scenarios cannot be assumed to determine it. The framework therefore reports the ranking from the full scenario set and tests its sensitivity to dropping single scenarios.

\section{Methodology: Robust Ranking and Validation}
\label{sec:g7-method}

\subsection{Metrics}
Let $M$ be the eligible cells of a country and $S$ the set of futures from the Latin hypercube design. For cell $i$ and future $s$, $C(i,s)$ is the LCOE. The regret uses the reference of Section~\ref{sec:g4-q}:
\begin{align}
R(i,s) &= \frac{C(i,s) - L^{*}_{s}}{L^{*}_{s}}, \qquad L^{*}_{s} = Q_{q}\!\big(C(\cdot,s)\big), \label{eq:g7-regret}\\
Z(i)   &= \max_{s \in S} R(i,s), \label{eq:g7-mmr}\\
\mathrm{Sat}(i;\tau) &= \frac{1}{|S|} \sum_{s \in S} \mathbf{1}\!\left[C(i,s) \le \tau\right]. \label{eq:g7-sat}
\end{align}
Cells are ranked by $Z(i)$ (lower is better) and by $\mathrm{Sat}(i;\tau)$ (higher is better) with $\tau$ and $q$ from Group~4. The top fraction $p_k \in \{5, 10, 20\}\%$ defines the robust set. The worst-case score $Z$ is sensitive to a few extreme futures, so the satisficing share is always reported with it, and the ranking is also reported after removing each future in turn.

\subsection{Scenario discovery}
PRIM is applied to the space of uncertain factors, and the unit is the future (a draw of the uncertain design combined with a climate scenario), not the cell. The nominal shortlist is the $k$ cells with the lowest nominal LCOE. A future is a case of interest when the share of the shortlist's capacity held by cells whose relative-regret rank in that future exceeds $2k$ reaches a threshold of 0.20. The box that maximizes the density of such futures subject to a minimum coverage is reported with its coverage and density. % transcribed from M-F7-08

\subsection{Validation against published potentials (F7b)}
For each country and technology, the modeled technical potential is compared with the bands of Table~\ref{tab:g7-pot} after harmonizing hub height, power density, exclusions and the DC or AC basis. The result is reported as a ratio to each band with its premises, and it is not a pass or fail test. The pipeline hub height is 100~m for all three countries, so the 100~m values are the primary comparators, and the 150~m layer is run once per country as a sensitivity. Capacity in operation from the Global Energy Monitor is compared with the modeled cell capacity as a second check. % verify: confirm GEM version in F7b
For Portugal no published area totals were found, so validation there rests on installed capacity and on the planning figure of Group~4, which is declared as a limitation.

\subsection{Spatial aggregation test}
The ranking is recomputed on a coarser aggregation of the 0.05$^{\circ}$ grid, and the overlap of the top-$k$ sets and the rank correlation between resolutions are reported. % verify: aggregation levels as defined in OQ-064

\subsection{Limitations specific to Group 7}
\begin{itemize}
\item Published potentials differ in hub height, density and exclusions, so ratios are indicative.
\item No area-based potential was found for Portugal.
\item The 143~GW value belongs to the 2001 atlas at 50~m and is not a value of the 2013 simulation edition, which publishes no national total in the sources read.
\item The reduction to $n+1$ scenarios is not available for discrete cell ranking; the full scenario set is used.
\end{itemize}
